\setlength{\absparsep}{18pt} % ajusta o espaçamento dos parágrafos do resumo
\begin{resumo}
	Este trabalho apresenta a reconstrução da plataforma \emph{web} do Pró-Mamá, programa municipal de aleitamento materno de Osório, sob a perspectiva da continuidade tecnológica e do apoio digital à saúde materno-infantil. A pesquisa caracteriza-se como aplicada, de natureza tecnológica e abordagem mista, organizada como estudo de caso único e apoiada por levantamento documental, análise de requisitos e desenvolvimento incremental. O trabalho revisou limitações da versão anterior, especialmente a ausência de comunicação segura, dificuldades de manutenção, baixa reprodutibilidade da infraestrutura e dependência de regras fixas no código. Como resultado, foram reestruturados a \emph{API}, o Painel Administrativo e a infraestrutura de implantação, com uso de tecnologias como Node.js, TypeScript, React.js, PostgreSQL, Redis, Docker, Nginx e Cloudflare. A solução passou a contar com autenticação e autorização por perfis, documentação automatizada por OpenAPI, processamento assíncrono de notificações, migração controlada de dados legados, publicação por HTTPS e rotinas de integração e entrega contínua. A validação automatizada concentrou-se no \emph{back-end}, totalizando 143 verificações aprovadas e cobertura global próxima de 80\% nas métricas de linhas e instruções. Conclui-se que a reconstrução entregou uma base mais segura, administrável e sustentável para a continuidade do Pró-Mamá, preservando seu valor social e ampliando sua capacidade de manutenção e expansão.
	
	\textbf{Palavras-chave}: aleitamento materno. Pró-Mamá. saúde digital. aplicações web. segurança da informação.
\end{resumo}
